% Transcription — fonds Alexandre Grothendieck, shelfmark 123, batch 2 (pages 21-40).
% Established from the facsimile archives/batches/123/batch-02.pdf (a mirror of
% the folder's PDF; no cover sheet, so PDF page k = folder page 20+k), per the
% skill transcribe-grothendieck.
%
% Pass: Opus 5.5 (claude-opus-5-5), 2026-09-24 — first pass, unchecked against the pages by a human.
%
% \dating{} carries the folder's own inventory date, "1969", verbatim; the
% letter on pages 21-22 is dated « Massy 22.10.1969 » in his hand. The group
% "119-121" (« Autour d'Esquisse d'un programme », 1968-1991) is not used.
%
% Pages. Sixteen of the twenty pages carry a \page{}. Absent:
%   - 23 and 36, covers bearing only a title in his hand (buff paper); the
%     titles head the two \section*s that follow them (hand.md §5):
%     23 « Diagrammes d'espaces homogènes liés aux éléments de Kostant, de
%     Coxeter et aux pts cycliques », 36 « Liste d'espaces homogènes
%     remarquables et leurs relations » (with a pencil continuation);
%   - 24 and 32, blank buff sheets.
%
% The letter (21-22). A typed letter in English, « Dear Kostant », signed
% A. Grothendieck, with his hand for a few insertions and a footnote. It is
% HIS letter to Kostant (the one announced by the cover of page 14, batch 1)
% and about the mathematics in hand, so it is transcribed in full, in
% English. The words he inked in are given as plain text and listed in a
% \note (they complete the typed sentence rather than correct it).
%
% The typed versos (27, 29, 31, 35, 39). The Kostant notes of pages 26, 28,
% 30, 34, 38 are written on the backs of typescript leaves of two other
% texts: an exposé on topos (§3.1, nos 3.1.2-3.1.3: exactness of u^*, u_*,
% U-topos; leaves 27 and 39, the first cancelled by a vertical stroke) and a
% paragraph « 2. Accouplements définis par une biextension » (nos 2.1-2.1.13;
% leaves 35, 29, 31, typed pages 6, 8, 11). Neither is named on the leaves,
% and neither bears on the folder's mathematics. They are not skipped
% because they carry marks in ink that look like his hand (hand.md §6: a
% typed verso is kept when it carries his hand). Following the brief's rule
% for typescripts (folders 100 and 71), each is SUMMARISED in one French
% \note, and only the handwritten marks are given: the substantive
% insertions as \marginal{}, the routine ones (inked Greek letters and
% asterisks, cross-reference numbers) named once in the note. This is a
% judgement call: they could equally be cut to nothing as unrelated versos
% (folder 5 batch 2, folder 31 batch 1). Whose hand did the correcting is
% not asserted beyond « semble de sa main ».
%
% What the batch is.
%   21-22  The letter to Kostant, Massy 22.10.1969: classify quadruples
%          (T,B,T',B'), T and T' maximal tori « in apposition », B ⊃ T and
%          B' ⊃ T' Borels; number of conjugacy classes card(W)^2/hzφ(h)
%          (footnote: N(T) ∩ N(T') extension of (Z/hZ)^* by Z/hZ × Z); the
%          stabiliser T ∩ T' = e makes G « rigid »; questions on B, B' in
%          general position, a distinguished class, automorphism groups;
%          then the rank of the tangent map of g → h/W (Kostant) and
%          G → T/W (Steinberg) on unipotent orbits of dimension n-r-2 (all
%          roots of the same length, Klein singularities), the dimension of
%          the variety of Borels containing an element, and rationality of
%          the singularities of the fibres (Brieskorn's depth argument).
%   25     A table of the notions and their abbreviations with, on the right,
%          a group (stabiliser): Kostant/Coxeter elements, subgroups,
%          couples, précouples (tori in apposition), cyclic lines, couples de
%          Kostant-Killing, normalised versions; « kif kif » in char. 0 once
%          ω ∈ μ_h^*(k) is chosen; margins in brown ink.
%   26     The lattice of the stabilisers (0, z, μ, μ×z, z·T, μ·T', N(μ),
%          N ∩ N', N') and a diagram of the spaces (Kost-Kill-norma down to
%          Tor, Toapp); definitions of C Kost norma and Kost Kill; footnotes
%          * and ** on when the isomorphisms are defined.
%   28     Couples of tori in apposition ≃ Coxeter subgroups μ; a 4×5 grid
%          of intersections of T, z(μ), N(μ), N, N(D), T', N' with their
%          quotients (« coc. » squares); the canonical extension of
%          (Z/hZ)^* by μ × z realised in W and W'.
%   30     The same grid for a Coxeter element γ: T'^γ ≃ P'/(1-c)P'(1) ≃ z;
%          to prove: every element of T normalising D lies in K(μ_h); a
%          determinant = Σn_i + 1; N(T) ∩ N(D) = N(μ) ∩ N(D).
%   33     « Quelques rappels de résultats de Kostant »: a) cyclic lines D
%          normalised by μ_h form a homogeneous space under T, in bijection
%          with the tori T' in apposition; b) cyclic lines mod μ_h^*
%          identify with a canonical line Δ, via Kostant's g → t/W;
%          g_ψ ≃ Δ (or Δ^{-1}, Δ^{⊗2} ≃ 1). Long oblique margin on
%          torsors and an obstruction in H^2(Q;E;μ_h). Much of the connecting
%          prose is fast and only partly read.
%   34     The scheme RigKos of « rigidifications de Kostant » (T,T',B,B'),
%          (h,p) = 1: smooth over S, G acting freely; RigKos/G finite étale
%          of rank W^2/φ(h)hz; the split case via an épinglage,
%          e = e_{α_1}+…+e_{α_r}+e_{-ψ}, T' its centraliser, (Q × Q')/Φ.
%   37     List of schemes: principal simple subgroups of rank 1 (SSP, with
%          épinglage, Killing couple, torus, Borel …), principal tori and
%          homomorphisms G_m → G, principal additive subgroups and regular
%          unipotents, E Cox / E Kost, cyclic lines, CTAP, CCox, CKost.
%   38     The diagram of these homogeneous spaces under G, and a diagram of
%          the corresponding groups (e, t, b_u, N_1, T, n, b, Z(b_u), N,
%          N_0, H, N(b_u)) with dimensions on the edges.
%   40     « Résultats de Kostant » on remarkable conjugacy classes in a
%          semisimple adjoint G over C: a) SL(2)-subgroups ≃ unipotent
%          classes ≃ admissible G_m → T with components 0, 1, 2; b) the
%          principal case, centralisers and normalisers; c) the principal
%          G_m → G evaluated at a primitive h-th root of 1 gives regular
%          elements whose image in W is a Coxeter element; d) eigenlines of
%          the Coxeter element on Lie(T). The page is crowded with
%          interlinear additions and several margins; many are fragmentary.
%
% Notation fixed for the batch, as in batch 1: \mathfrak{g}, \mathfrak{t},
% \mathfrak{h} for his gothic letters; \mathfrak{z} for the curly 3 he uses
% for the centre (and its order z); \mu, \mu_h; W, W' the Weyl groups of T,
% T'; the abbreviations (E Kost, SSPEP, CTAP, …) are kept as he wrote them.
% Headless strokes in his diagrams (inclusions, not maps) are `no head`.
\documentclass[11pt,a4paper]{article}
% Le préambule commun à toutes les transcriptions du fonds Grothendieck.
%
% Il tient en une idée : **l'incertitude se compose, elle ne se résout pas.**
% Une transcription qui lisse un mot douteux a détruit la seule chose pour
% laquelle on la faisait — un lecteur doit pouvoir savoir, à chaque endroit,
% ce qui était sur le feuillet et ce qui a été ajouté. Les six macros qui
% suivent sont donc l'appareil critique, et non de la décoration.
%
% Les mêmes commandes sont reconnues par `scripts/render.mjs`, qui produit la
% vue de lecture du site. Une macro ajoutée ici doit l'être aussi là-bas,
% faute de quoi le rendu échouera bruyamment — ce qui est voulu : un rendu
% silencieusement faux est pire qu'un rendu absent.

\ProvidesPackage{grothendieck}[2026/08/08 Transcriptions du fonds Grothendieck]

\usepackage{amsmath,amssymb,amsthm}
\usepackage{mathrsfs}
\usepackage{stmaryrd}
% Les diagrammes commutatifs sont partout dans ce fonds : c'est la langue
% même de la géométrie algébrique de Grothendieck, et une transcription qui
% les rendrait en prose ne transcrirait rien.
\usepackage{tikz-cd}
% Français par défaut — c'est la langue des feuillets — mais l'anglais est
% chargé aussi : la traduction et le résumé sont des documents anglais, et la
% typographie française (espaces avant les deux-points) y serait une faute.
% Ces documents basculent par \selectlanguage{english} en tête de corps.
\usepackage[english,french]{babel}
\usepackage{fontspec}
\usepackage{geometry}
\geometry{a4paper,margin=25mm,marginparwidth=28mm}
\usepackage{marginnote}
\usepackage{xcolor}
% ulem, not soul. `soul`'s \st and \ul are built on a character-by-character
% scan that XeLaTeX's font machinery breaks: the document fails with an
% undefined control sequence rather than degrading. ulem does the same two jobs
% and is Unicode-engine safe. `normalem` stops it hijacking \emph, which the
% transcriptions use for Grothendieck's own emphasis.
\usepackage[normalem]{ulem}
\usepackage{hyperref}
\hypersetup{colorlinks=true,linkcolor=black,urlcolor=black,pdfborder={0 0 0}}

% --- Métadonnées du lot -----------------------------------------------------
% Reprises telles quelles par le script de rendu, qui les affiche en tête de
% la vue de lecture. Elles fixent ce que le fichier prétend couvrir : sans
% elles, une transcription flotte sans point d'attache dans un fonds de seize
% mille feuillets.

\newcommand{\folder}[1]{\gdef\grothFolder{#1}}
\newcommand{\batch}[1]{\gdef\grothBatch{#1}}
\newcommand{\leaves}[2]{\gdef\grothFirst{#1}\gdef\grothLast{#2}}
\newcommand{\foldertitle}[1]{\gdef\grothFoldertitle{#1}}
\gdef\grothFolder{?}\gdef\grothBatch{?}\gdef\grothFirst{?}\gdef\grothLast{?}\gdef\grothFoldertitle{}

% --- Le résumé --------------------------------------------------------------
%
% Il ouvre la lecture modernisée, sous le titre « Résumé », et il est composé
% en petit. Ce n'est pas un ornement : c'est le seul passage du document qui
% ne soit pas des mathématiques, et un lecteur doit pouvoir le reconnaître
% avant de l'avoir lu — puis le sauter s'il connaît déjà le sujet, ou s'y
% arrêter s'il ne le connaît pas. Le filet qui le suit marque la reprise du
% corps.

\newenvironment{resume}
  {\small\setlength{\parskip}{0.45em}}
  {\par\medskip\hrule\medskip}

% --- Mots-clés ---------------------------------------------------------------
%
% En anglais, en clôture du résumé de la lecture modernisée : le vocabulaire
% moderne sous lequel chercher ce que les feuillets construisent. Le manifeste
% du site (`npm run manifest`) les extrait pour étiqueter la cote entière —
% cette ligne est la source unique des tags, il n'y a pas de fichier à côté.
\newcommand{\keywords}[1]{\par\smallskip\noindent
  {\footnotesize\textsc{Keywords} --- \textit{#1}}\par}

% --- L'appareil critique ----------------------------------------------------

% Le numéro de feuillet, en marge. C'est l'ancre qui permet de régler un
% désaccord sur une formule en regardant *un* feuillet plutôt qu'une chemise
% de six cents. Le site s'en sert aussi pour tourner les pages du fac-similé
% quand on fait défiler la transcription.
\newcommand{\leaf}[1]{%
  \marginnote{\footnotesize\color{gray!70}\textsf{#1}}%
  \label{leaf:#1}\ignorespaces}

% Illisible : on ne devine pas. Un mot inventé qui se lit comme les autres est
% le pire résultat possible.
\newcommand{\ill}{\textcolor{red!60!black}{[\,\ldots\,]}}

% Lecture douteuse : le mot est donné, et signalé comme incertain.
\newcommand{\uncertain}[1]{\uline{#1}}

% Ajout de l'éditeur — une lettre manquante, un mot sous-entendu.
\newcommand{\add}[1]{\textcolor{blue!50!black}{[#1]}}

% Biffé par Grothendieck lui-même. Ce qu'il a rayé fait partie du document :
% une rature dit souvent où la pensée a changé de direction.
\newcommand{\struck}[1]{\sout{#1}}

% Note du transcripteur, sur l'état matériel du feuillet.
\newcommand{\note}[1]{\par\smallskip{\small\color{gray!60!black}\textit{[#1]}}\par\smallskip}

% Note marginale de Grothendieck, distincte du corps du texte.
\newcommand{\marginal}[1]{\par\smallskip\noindent\hspace{1em}%
  {\small\color{gray!60!black}#1}\par\smallskip}

% --- Titre ------------------------------------------------------------------

\AtBeginDocument{%
  \begin{center}
    {\large\bfseries Fonds Alexandre Grothendieck --- cote n\textsuperscript{o}~\grothFolder}\\[2pt]
    {\itshape\grothFoldertitle}\\[4pt]
    {\small feuillets \grothFirst--\grothLast\ (lot \grothBatch)}\\[2pt]
    {\footnotesize\color{gray!60!black}Transcription établie d'après le fac-similé numérisé
     par l'Université de Montpellier.\\
     Les lectures incertaines sont soulignées, les illisibles notées [\,\ldots\,],
     les ajouts entre crochets.}
  \end{center}
  \vspace{1em}}

\endinput


\folder{123}
\batch{2}
\pages{21}{40}
\dating{1969}
\watermark{Édition de démonstration}
\foldertitle{[Autour de Kostant] : tapuscrits (s.d.), notes manuscrites (s.d.), lettre (1969).}

\begin{document}

\section*{Lettre à Kostant}

\page{21}
\note{le titre est de l'éditeur. Lettre dactylographiée, en anglais, de sa main pour le lieu, la date,
quelques corrections, un appel de note et la note elle-même au pied du
feuillet ; signée de sa main à la page suivante. Les fautes de frappe corrigées
à la machine ne sont pas relevées.}

Massy 22.10.1969

Dear Kostant,

I read again through your papers in the Amer.\ Journ.\ on « The principal
three-dimensional subgroup … » and « Lie group representations on
polynomial rings », prompted by some nice work of Brieskorn on Klein
singularities. I very much appreciate this work of yours, and would appreciate
getting reprints of any later work you may have available. As I already felt
when your work appeared, it cries for careful reconsideration in the framework
of Alg.\ Geometry, and suggests number of interesting problems. I wonder if you
know the answer to some of them. For instance, I feel it would be of interest to
classify (for a given complex simple adjoint Lie group, say) quadruples
$(T,B,T',B')$, where $T$ and $T'$ are maximal tori « in apposition » (in your
terminology), and $B$ and $B'$ Borel subgroups of $G$ containing $T$ resp.\
$T'$. The number of conjugacy classes of such animals is equal to
$\operatorname{card}(W)^{2}/hz\varphi(h)$, where $W$ is the Weyl group,
$h$ the Coxeter number, $z$ the order of the center of $\tilde{G}$, and
$\varphi(h) = \operatorname{card}(\mathbb{Z}/h\mathbb{Z})^{*}$ is the Euler
indicatrix; if you want to classify such data with moreover a
generating element (principal regular in your terminology) given for the
finite group of order $h$ $T \cap N(T')$, then we get
$\operatorname{card}(W)^{2}/hz$ choices. One of the reasons I think this
question is interesting is that the stablity\note{sic} subgroup in
$G$ of such a quadruple is $T \cap T' = e$,$^{*}$ hence such a structure makes
$G$ entirely « rigid » up to exterior automorphisms. Among the first questions
one may wish to ask in this direction are the following: is it true that for
some, or for all, such quadruples, $B$ and $B'$ are « in general position »
i.e.\ $B \cap B'$ is a torus (necessarily maximal), so that $B$ and $B'$ are
« opposite » to each other with respect to the latter ? Is there in any sense a
\emph{distinguished} conjugacy class of such quadruples (which should then be,
of course, invariant by action of exterior automorphisms) ? What are the groups
of automorphisms of such quadruples (they
\note{les deux exposants $2$ de $\operatorname{card}(W)$, « of $\tilde{G}$ » et
« generating » sont ajoutés à la main ; « distinguished » est souligné}

\marginal{$^{*}$ NB : $N(T) \cap N(T')$ is an extension of
$(\mathbb{Z}/h\mathbb{Z})^{*}$ by $\mathbb{Z}/h\mathbb{Z} \times Z$
\struck{\ill{}}, hence of order $hz\varphi(h)$}
\note{note de sa main au pied du feuillet ; sous $Z$, relié par un trait :
« center of $\tilde{G}$ »}

\page{22}
are isomorphic to subgroups of the group of exterior automorphisms of $G$) ?

In a different direction, it would be interesting to know more about the
canonical morphism $\mathfrak{g} \to \mathfrak{h}/W$, where $\mathfrak{g}$ is
the Lie algebra of $G$ and $\mathfrak{h}$ a Cartan subalgebra, and of the orbits
of $G$ on $\mathfrak{g}$, and the analoguous morphism of Steinberg
$G \to T/W$ and the orbits of $G$ acting on itself. Do you know for instance
exactly how varies the rank ($=$ rang of the tangent map) of this map ? There
are some reasons to believe that for points of unipotent orbits of dimension
$n-r-2$ (just the next lower after the maximal dimension $n-r$), the
rank is $r-1$ in case all roots have same length (that is the corresponding
diagrams $A_{\uncertain{n}}, D_{\uncertain{n}}, E_6, E_7, E_8$ are those
corresponding to Klein singularities), and $< r-1$ in all other cases; in
the first case, I would expect that there is just one orbit as stated, by the
way, and maybe this property will be caracteristic of the « homogeneous » case
when all roots have same length. Also, what can be said of the dimension of the
set of all Borel subgroups containing a given element $g \in G$, resp.\ whose
Lie algebra contains a given element $x \in \mathfrak{g}$ ? If $2N = n-r$ is
the number of roots and if the dimension of the orbit of the given element is
$2N-2i$, I would expect the dimension of the previous variety to be equal to
$i$; the inequality $\leqslant i$ would imply, by generalizing a depth argument
of Brieskorn, that the singularities of the fibers of your map
$\mathfrak{g} \to \mathfrak{h}/W$ (resp.\ Steinberg's $G \to T/W$ in case $G$
is simply connected instead of adjoint) are « rational », i.e.\ if $F$ is such a
fiber and $F' \xrightarrow{f} F$ a resolution of singularities of $F$, then
$R^{i} f_{*}(\mathcal{O}_{F'}) = 0$ for $i > 0$. Maybe the answer to most of
the questions I am asking are well known and I am just ignorant; for some of
them, it would indeed be scandalous that the experts do not know the answer !
In any case, I would be grateful for any comment you would have on any of my
questions.

Sincerely yours

A.~Grothendieck
\note{« $n-$ » et « all » sont ajoutés à la main ; l'indice des diagrammes
$A$ et $D$ est frappé par-dessus un autre caractère et se lit mal ; la
signature est manuscrite}

\section*{Diagrammes d'espaces homogènes liés aux éléments de Kostant, de Coxeter et aux pts cycliques}

\note{titre de sa main, en haut à droite d'un feuillet de papier bis qui ne
porte rien d'autre (page 23) : le feuillet est une chemise et ne reçoit donc
pas de \texttt{page}}

\page{25}
\note{tableau en trois colonnes : à gauche les notions, au milieu leurs sigles,
à droite un groupe (le stabilisateur, semble-t-il, de la donnée correspondante) ;
on le donne ligne par ligne}

\begin{itemize}
  \item éléments de Kostant $\kappa : \mu_h \to G$ ; éléments de Coxeter
  $\gamma \in G(k)$. --- E Kost, E Cox : kif-kif si car.~$0$ et si on se donne
  $\omega \in \mu_h^{*}(k)$. ---
  \struck{ext.\ de $\mathfrak{z}$ par $T$} \struck{\ill{}}
  $\mathfrak{z}.T$ ($\tfrac12$ direct)
  \item \struck{Ss-groupes} Ss-groupes de Kostant $\mu \hookrightarrow G$ ;
  ss-groupes de Coxeter $\mu \hookrightarrow G$. --- SG Kost, SG Cox : kif kif
  en car.~$0$. --- ext.\ de $(\mathbb{Z}/h\mathbb{Z})^{*}$ par
  $\mathfrak{z}.\struck{\ill{}}T$
  \item Couple de Kostant $(\kappa : \mu_h \to G, T')$ [ou on se donne
  $(\kappa : \mu \hookrightarrow G, D)$, $D \in \hat{\mathfrak{g}}(1)^{\mathrm{reg}}$] ;
  couple de Coxeter $(\gamma, T)$. --- C Kost, C Cox : kif kif si car.~$0$ et
  si on se donne $\omega \in \mu_h^{*}(k)$. --- $\mu_h \times \mathfrak{z}$
  \item précouple de Kostant $(\mu \subset G, T')$ ; précouple de Coxeter
  $(\mu \subset G, T)$ ; \uncertain{i.e.}\ couple de tores en apposition. ---
  \struck{PC Kost, PC Cox : kif kif en car.~$0$} Toapp$^{*}$. --- ext.\ de
  $(\mathbb{Z}/h\mathbb{Z})^{*}$ par $\mu_h \times \mathfrak{z}$, (ou produit)
  \item droites cycliques $D$ ; \struck{Classe} \struck{\ill{}} de Coxeter
  $(T, \dot{\gamma})$, $\dot{\gamma} \in \hat{W}(h)$ ; élément
  \struck{cyclique} normalisé (dépend du $\delta$ Killing). --- D Cy,
  \struck{\ill{}} Cox, E Cy$_a$ : kif kif si car.~$0$ et on se donne
  $\omega \in \mu_h^{*}(k)$. --- $\mu_h.T'$ ($\tfrac12$ direct) ; $T'$
  \item couple \struck{triple} de Kostant-Killing
  $((B,T),T') \sim (B,T,D)$ (tel que $T'$ normalise un
  $\mu_h \xrightarrow{\kappa} T$ défini par $B \supset T$). ---
  \struck{\ill{}} Kost Kill. --- $\mu_h$
  \item couple \struck{triple} de \struck{Coxeter} Kostant normalisé
  $(T, \struck{\ill{}}x)$ \struck{$\gamma$ cyclique normalisé}
  \struck{droite cyclique \ill{}} un tore $T'$ en apposition avec $T$, i.e.
  \ill{}. --- \struck{PCox} C Kost norma. --- $\mathfrak{z}$
  \item Couple de Kostant-Killing normalisé $((T,B),x)$. --- Kost Kill norma.
  --- $1$
\end{itemize}
\note{les mots « couple » au-dessus de « triple » biffé, « Killing » après
« Kostant » à la sixième ligne et « Kostant normalisé » au-dessus de
« Coxeter » biffé à la septième sont des surcharges ; les sigles ont été
récrits plusieurs fois et les biffures y sont lourdes}

\marginal{\uncertain{Vrais} seulement en car.~$0$ (ou car.\ $p > h$ ?)}
\note{à l'encre brune, dans la marge gauche, en face de l'accolade qui réunit
les deux couples de la troisième ligne ; en face de la quatrième, de la même
encre : « id », qui se lit mal}

\marginal{$\delta$ choix d'un $\alpha \in \Delta^{*}(k)$, où $\Delta$ est le
groupe commun aux droites cycliques par $\mu_h$ (\emph{NB} modulo signe, il y a
un choix privilégié \struck{\ill{}} d'un $x \in \mathfrak{z}$, \uncertain{des
racines} \uncertain{par} coef.\ \uncertain{géodésiques})}
\note{à l'encre brune, au pied du feuillet, sous un trait : c'est la note
appelée par le $\delta$ de la cinquième ligne}

\page{26}
\note{sur le verso d'un feuillet dactylographié, dont le texte transparaît à
l'envers. En haut, un treillis de sous-groupes, les traits sans flèche
marquant l'inclusion ; les étiquettes sont celles qu'il porte sur les traits}

\begin{tikzcd}[column sep=small, row sep=small, nodes={font=\scriptsize}]
  & & 0 \arrow[dl, no head] \arrow[dr, no head] & & \\
  & \mathfrak{z} \arrow[dr, no head] & & \mu \arrow[dl, no head] & \\
  & & \mu \times \mathfrak{z} \arrow[dl, no head, "T/\mu"'] \arrow[dd, no head, "(\mathbb{Z}/h\mathbb{Z})^{*}"] \arrow[dr, no head] & & \\
  & \mathfrak{z}\cdot T \arrow[dl, no head, "(\mathbb{Z}/h\mathbb{Z})^{*}"'] & & \mu \cdot T' \arrow[dr, no head] & \\
  N(\mu) \arrow[rr, no head] \arrow[d, no head] & & N \cap N' \arrow[rr, no head] & & N' \\
  N & & & &
\end{tikzcd}

En dessous, un second diagramme, d'espaces cette fois :

\begin{tikzcd}[column sep=small, row sep=small, nodes={font=\scriptsize}]
  & & & \text{Kost-Kill-norma} \arrow[dl, no head] \arrow[dr, no head] & & & \\
  & & \text{C Kost norma} \arrow[dr, no head] & & \text{Kost Kill} \arrow[dl, no head] & & \\
  & & & \text{C Kos} \arrow[dl, no head] \arrow[ddd, no head] \arrow[dr, no head, "\wr^{*}"] \arrow[drr, no head] & & & \\
  & & \text{E Kos} \arrow[dl, no head] \arrow[d, no head, "\wr^{**}"] & & \text{C Cox} \arrow[dll, no head] \arrow[ddl, no head, "{}^{*}"] \arrow[dr, no head] & \text{D Cy} \arrow[d, no head, "\wr^{**}"] \arrow[dr, no head] & \\
  & \text{SG Kos} \arrow[dl, no head] \arrow[dr, no head, "\wr^{*}"] & \text{E Cox} \arrow[d, no head] & & & \text{P Cox} \arrow[dr, no head] & \text{Tor} \arrow[d, no head, "="] \\
  \text{Tor} \arrow[dr, no head, "="'] & & \text{SG Cox} \arrow[r, no head, "{}^{*}"] & \text{Toapp} \arrow[rrr, no head, "{}^{*}"'] & & & \text{Tor} \\
  & \text{Tor} \arrow[ur, no head] & & & & &
\end{tikzcd}

\note{les « $=$ » rendent deux doubles traits verticaux « $\parallel$ » entre
deux « Tor » ; le long trait qui descend de C Kos passe derrière C Cox et
s'arrête à Toapp, lecture incertaine. Au-dessus de C Kost norma et de Kost
Kill, sous « Kost-Kill-norma », une ligne encadrée et biffée :
\struck{$((T,B),(T',B'))$ \ill{} $T, T'$ en apposition \ill{}}. Sous les deux
sommets de la deuxième ligne, leurs définitions, que les traits du diagramme
traversent :}

C Kost norma : paires $(T,(T',B'))$, avec $T$ en apposition avec $T'$, $B'$
dans une classe donnée \uncertain{mod} $W'$ (ss-gr.\ extension de
$(\mathbb{Z}/h\mathbb{Z})^{*}$ par $\mu$).

Kost Kill : paires $((T,B),T')$, telles que $(\kappa, T')$ soit une paire de
Kost., $\kappa$ étant l'élément de Kostant défini par le couple de Killing.

$^{**}$ Isom défini si on travaille \uncertain{sur} $\mathring{\mu}_h^{*}$.

$^{*}$ Isom défini si $(h,p) = 1$ ou notion définie dans ce cas.

\note{une grande courbe au crayon enferme la partie haute du second diagramme
(de Kost-Kill-norma à C Cox) ; un trait en descend jusqu'à :}
leur existence dans cas quasi-déployé \uncertain{suit} par descente.

\page{27}
\note{Feuillet de tapuscrit, paginé « 72 », barré d'un long trait vertical et
remployé comme brouillon : un exposé sur les topos, §~3.1 (n°~3.1.2) — si
$u \colon E \to E'$ est un morphisme de topos, $u_*$ commute aux limites
projectives et $u^*$ aux limites inductives et aux limites projectives finies,
donc $u^*$ est exact ; pour toute espèce de structure algébrique $\Sigma$
définie par des limites projectives finies et des limites inductives
quelconques, l'image inverse par $u^*$ d'un objet de $E'$ muni d'une structure
d'espèce $\Sigma$ en est muni ; exemples : groupe, anneau, module, comodule,
bigèbre, torseur. Le texte ne touche pas aux notes sur Kostant. Les corrections
à l'encre bleue, qui semblent de sa main, sont d'un correcteur de tapuscrit :
astérisques de $u_*$, $u^*$ et lettre $\Sigma$ encrées, renvois
« (I~2.11) », « (3.2.1) », « (1.5 et 1.8) », « $\mathcal{U}$- » devant
« limites projectives » ; on ne relève que les ajouts de mots :}

\marginal{est exact à gauche i.e.}
\note{au-dessus de « commute aux limites projectives finies », qu'il remplace}

\marginal{On peut donc dire que}
\note{en tête de « c'est le foncteur image inverse $u^*$ qui, dans le couple
$(u_*,u^*)$, possède les propriétés d'exactitude les plus remarquables »,
où « donc » est biffé}

\marginal{et pour}
\note{devant « tout objet de $E'$ muni d'une structure d'espèce $\Sigma$ »}

\marginal{l'expliciter et de}
\note{inséré dans « nous conseillons au lecteur de se convaincre de sa
validité » ; à la note de bas de page, « vérifier » est biffé, et la fin
« c'est une conséquence immédiate de 1.2 » est barrée}

\page{28}
Couples de tores $(T,T')$ en apposition

$\simeq$ ss-groupes de Coxeter $\mu$, \struck{\ill{}} de Cartan
\uncertain{normalisé} \uncertain{tel que} $\mu$ opère \uncertain{fidèlement}
sur une droite \uncertain{cyclique} de Cartan ($\mathrm{Lie}\,\mathfrak{g}$)

\uncertain{et} le choix d'un \uncertain{isom.}\ $\mu \simeq \mu_h$ équivaut au
choix d'un $D$ …
\note{ces trois lignes sont encadrées ; à droite, hors du cadre, une accolade
réunit « sous-groupes de Kostant-Coxeter » et « Éléments de Kostant », et en
dessous « \uncertain{Sous-g} »}

\note{puis un tableau de sous-groupes, en quatre lignes reliées par des traits
verticaux ; sur les traits horizontaux et verticaux, les quotients ; « coc. »
(cocartésien, semble-t-il) marque chaque carré. On le redonne tel qu'il se
lit :}

\begin{tikzcd}[column sep=small, row sep=small, nodes={font=\scriptsize}]
  T \arrow[r, no head, "\mathfrak{z}"] \arrow[d, no head, "T/\mu"'] & \mathfrak{z}(\mu) \arrow[r, no head, "(\mathbb{Z}/h\mathbb{Z})^{*}"] \arrow[d, no head, "T/\mu"] & N(\mu) \arrow[r, no head] \arrow[d, no head, "T/\mu"] & N \arrow[d, no head] & \\
  T \cap N' \arrow[r, no head, "\mathfrak{z}"] \arrow[d, no head, "0"'] & \mathfrak{z}(\mu) \cap N' \arrow[r, no head, "(\mathbb{Z}/h\mathbb{Z})^{*}"] \arrow[d, no head, "0"] & N(\mu) \cap N' \arrow[r, no head, "="] \arrow[d, no head, "(\mathbb{Z}/h\mathbb{Z})^{*}"] & N \cap N' \arrow[r, no head] \arrow[d, no head, "(\mathbb{Z}/h\mathbb{Z})^{*}"] & N' \arrow[d, no head] \\
  T \cap N(D) \arrow[r, no head, "\mathfrak{z}"] \arrow[d, no head, "\mu"'] & \mathfrak{z}(\mu) \cap N(D) \arrow[r, no head, "0"] \arrow[d, no head, "\mu"] & N(\mu) \cap N(D) \arrow[r, no head, "="] \arrow[d, no head, "\mu"] & N \cap N(D) \arrow[r, no head, "T'/\mathfrak{z}"] \arrow[d, no head, "\mu"] & N(D) = \mu \cdot T' \arrow[d, no head, "\mu"] \\
  T \cap T' \arrow[r, no head, "\mathfrak{z}"] & \mathfrak{z}(\mu) \cap T' \arrow[r, no head, "="] & N(\mu) \cap T' \arrow[r, no head, "="] & N \cap T' \arrow[r, no head, "T'/\mathfrak{z}"] & T'
\end{tikzcd}

\note{le signe $\mathfrak{z}$ rend la lettre en forme de 3 qu'il emploie ici
pour le centre ; sous $T \cap T'$, « $0$ » avec des guillemets de répétition ;
au-dessus de $N \cap N(D)$, entre $=$ et le trait suivant, « $\mu \times
\mathfrak{z}$ » ; sous la dernière flèche horizontale, une étiquette biffée,
illisible}

\struck{\ill{}}
\[
\begin{cases}
  N \cap N(D) = \mathfrak{z}(\mu) \cap N(D) \\
  N \cap T' = \mathfrak{z}(\mu) \cap T' \\
  N \cap N' = N(\mu) \cap N'
\end{cases}
\]

\begin{tikzcd}[column sep=small, row sep=small, nodes={font=\scriptsize}]
  T \arrow[r, no head] \arrow[d, no head] & T \cdot \mathfrak{z} \arrow[r, no head, "(\mathbb{Z}/h\mathbb{Z})^{*}"] \arrow[d, no head] & N(\mu) \arrow[r, no head] \arrow[d, no head] & N \arrow[d, no head] & \\
  \mu \arrow[r, no head] \arrow[d, no head] & \mu \times \mathfrak{z} \arrow[r, no head, "(\mathbb{Z}/h\mathbb{Z})^{*}"'] \arrow[d, no head] & N(\mu) \cap N' \arrow[r, no head, "="] & N \cap N' \arrow[r, no head] & N' \\
  0 \arrow[r, no head] & \mathfrak{z} & & &
\end{tikzcd}

On a une extension \uncertain{can.}\ de $(\mathbb{Z}/h\mathbb{Z})^{*}$ par
$\mu \times \mathfrak{z}$, les deux extensions \uncertain{induites} de
$(\mathbb{Z}/h\mathbb{Z})^{*}$ par $\mu$ et de $(\mathbb{Z}/h\mathbb{Z})^{*}$
par $\mathfrak{z}$ se \uncertain{réalisent} \uncertain{canoniquement} dans
$W'$ et dans $W$. \uncertain{Cela} \uncertain{fait} \uncertain{dans} \ill{} un
groupe de Weyl des sous-groupes \uncertain{remarquables}, \uncertain{définis}
\uncertain{à} \uncertain{conjugaison} \ill{} …

\page{29}
\note{Feuillet de tapuscrit, paginé « 8 », verso du précédent ; d'un autre
texte que celui de la page 27 : un paragraphe « Accouplements définis par une
biextension », n°~2.1.8 à 2.1.10 — réduction au cas
$P = Q = (\mathbb{Z}/n\mathbb{Z})_T$ ; l'homomorphisme
$\alpha \colon (\mathbb{Z}/n\mathbb{Z})_T \otimes (\mathbb{Z}/n\mathbb{Z})_T
\to \mathrm{Tor}_1^{\mathbb{Z}}((\mathbb{Z}/n\mathbb{Z})_T,
(\mathbb{Z}/n\mathbb{Z})_T)$ est l'isomorphisme (2.1.8.1) déduit de la
résolution plate canonique ; définition symétrique
$\beta^n_{P,Q} \colon {}_nP \otimes {}_nQ \to \mathrm{Tor}_1(P,Q)$
(2.1.9.1)-(2.1.9.2), et la relation (2.1.9.3) $\beta^n_{P,Q} =
-\alpha^n_{P,Q}$ ; Lemme 2.1.10 : un certain diagramme est
anticommutatif. Sans rapport avec les notes sur Kostant. Les lettres grecques
et les exposants $n$ y sont encrés à la main ; le numéro « (2.1.8.2) » est
ajouté ; dans la marge, entouré : « $(\mathbb{Z}/n\mathbb{Z})_T$ », qui
remplace « $P = (\mathbb{Z}/n\mathbb{Z})_T$ » biffé dans la ligne ;
« premier » est souligné}

\page{30}
\marginal{Car si $g \in \mathfrak{z}(\gamma)$, $g \cdot D = D'$ est
\uncertain{conjugué} de $D$ par \uncertain{élément} de $T$, donc \uncertain{on
peut} \uncertain{modifier} $g$ \uncertain{pour} \uncertain{qu'il} \uncertain{soit}
dans $\mathfrak{z}(\gamma) \cap T'$}
\note{en tête du feuillet, relié par un trait au diagramme qui suit}

\begin{tikzcd}[column sep=small, row sep=small, nodes={font=\scriptsize}]
  T \arrow[r, no head, "\mathfrak{z}"] \arrow[d, no head] & \mathfrak{z}(\gamma) \arrow[r, no head] \arrow[d, no head] & N(T) \arrow[r, no head] \arrow[d, no head] & G \arrow[d, no head] \\
  T \cap N(T') \arrow[r, no head, "\mathfrak{z}"] \arrow[d, no head, "\mathbb{Z}/h\mathbb{Z}"] & \mathfrak{z}(\gamma) \cap N(T') = N'^{\gamma} \arrow[r, no head] \arrow[d, no head, "\mathbb{Z}/h\mathbb{Z}"] & N(T) \cap N(T') \arrow[r, no head] \arrow[d, no head] & N(T') \arrow[d, no head, "W'"] \\
  T \cap T' \arrow[r, no head, "\mathfrak{z}"] & \mathfrak{z}(\gamma) \cap T' = T'^{\gamma} \arrow[r, no head, "="] & N(T) \cap T' \arrow[r, no head] & T'
\end{tikzcd}

\note{devant $T \cap N(T')$, « $\gamma \in$ ». Sous le diagramme :}
$T \cap T' = 0$, et
$T'^{\gamma} \simeq P'/(1-c)P'(1) \simeq \mathfrak{z}$.
\note{une flèche remonte de $0$ vers la phrase suivante, une grande courbe
enveloppe la colonne de $\mathfrak{z}(\gamma)$}

Car tout élément de $T$ qui normalise $D$ est \struck{\ill{}} l'élément
neutre.

Car tout élément de $W$ qui centralise $c$ est une puissance de~$c$.

\emph{À prouver} : \uncertain{Tout} élément de $T$ qui normalise $D$ est dans
$K(\mu_h)$, i.e.\ $h\psi$ (\ill{}, \ill{} choix, les \ill{}) \uncertain{sont}
\uncertain{complétées} \uncertain{à} coeff.\ \uncertain{entiers}
\uncertain{définissant} $\alpha_i + \psi$.

\uncertain{Prouvons} \uncertain{donc} \ill{} \ill{}
\[
  \det \begin{pmatrix}
    n_1 + 1 & n_2 & n_3 & \cdots & n_r \\
    n_1 & n_2 + 1 & n_3 & \cdots & n_r \\
    \vdots & & & & \\
    n_1 & n_2 & & \cdots & n_r + 1
  \end{pmatrix} = \textstyle\sum n_i + 1
\]
\note{la matrice est esquissée : la première colonne de la deuxième ligne
porte « $n\,x$ », la dernière ligne commence par un $n_1$ surchargé}

Un élément de $N(T) \cap N(D)$ \struck{\ill{}} normalise $D$, et normalise
$T$, donc il \struck{\ill{}} normalise \struck{\ill{}}
$\mu_h = \mathbb{Z}/h\mathbb{Z} \subset T$ qui est formé des éléments de $T$
qui normalisent $D$. Donc
\[
  N(T) \cap N(D) = \struck{T \cap \mathrm{Norm}(\mu)}\ N(\mu) \cap N(D)
  \qquad (\mu \subset T \text{ par } \gamma), \quad
  \mathbb{Z}/h\mathbb{Z} \simeq \mu_h .
\]
\note{entre « $\subset T$ » et « par $\gamma$ », un mot biffé, illisible}

\page{31}
\note{Même tapuscrit, feuillet paginé « 11 » : fin de la démonstration de
2.1.11 (un triangle $P/nP \to {}_nP$, ${}_{n'}P$ par $x \mapsto mx$ et
$x \mapsto x$), Corollaire 2.1.12 (commutativité du carré (2.1.12.1) entre
$\alpha^{n'}_{P,Q}$ et $\alpha^{n}_{P,Q}$, les flèches horizontales induites
par $m \cdot \mathrm{id}$), puis 2.1.13 : pour $\ell$ premier,
$T_\ell(E) = ({}_{\ell^\nu}E)_{\nu \geqslant 0}$, et les $\alpha$ définissent un
homomorphisme de systèmes projectifs (2.1.13.2)
$\alpha^{(\ell)}_{P,Q} \colon T_\ell(P) \otimes T_\ell(Q) \to
T_\ell(\mathrm{Tor}_1(P,Q))$. Sans rapport avec les notes sur Kostant. Les
lettres grecques, les exposants et « $m \otimes m$ », « $m$ » sur les flèches
sont encrés à la main ; « qui montre que le diagramme envisagé » est biffé et
« évidemment » ajouté devant « commutatif » ; au bas, dans la marge,
« $\alpha^{\ell^\nu}_{P,Q}$ », et sous « pour $\nu$ variable » :}
\marginal{fonctoriel en $P$, $Q$ :}

\page{33}
Quelques rappels de résultats de Kostant :

a) Pour $\mu_h \hookrightarrow T \hookrightarrow G$ fixé, les droites
cycliques $D$ normalisées par $\mu_h$ \uncertain{(i.e.\ \ill{})}

1) forment un espace homogène sous $T$, stabilisateur $\mu_h$, donc un espace
homogène sous $\mathfrak{t}$. \uncertain{Ces} droites cycliques $D$
correspondent biunivoquement aux sous-tores $T'$ de $G$ \struck{\ill{}}
\uncertain{en apposition} \struck{$\mu_h$} à $T$ par $\mu_h$ (i.e.\ normalisés
par $\mu_h$, et tels que $\mu_h$ y admette une \struck{\ill{}} car.\ propre au
moins qui est \struck{\ill{}} car.\ primitif, $\Leftrightarrow$ la car.~$1$ y
est car.\ propre) : à $D$ correspond $T' =$ \uncertain{Cent}$(D)$, à $T'$
correspond \struck{l'unique droite des \ill{}} \uncertain{l'unique}
\ill{} \ill{} \uncertain{propre} de $\mu_h$ opérant sur $\mathrm{Lie}\,T'
\subset \mathfrak{g}$. On a d'ailleurs
$T \cap N(T') = \mu_h$.

b) Les droites cycliques $D$ de $\mathfrak{g}$ \emph{mod} $\mu_h^{*}$
s'identifient can., \uncertain{via} une droite $\Delta$ canonique sur $k$. On
l'obtient par descente à partir de $D^{\otimes h}$ défini sur le schéma des
droites cycliques \uncertain{ou} \uncertain{sur} \uncertain{le} schéma des
droites sur le schéma des tores maximaux $T'$, et à ce titre se descend. On
l'obtient aussi comme l'image des lieux cycliques de $\mathfrak{g}$ par
l'hom.\ de Kostant $\mathfrak{g} \to I = \mathfrak{t}/W$, qui a une structure
\emph{vectorielle} \uncertain{car} le stabilisateur (dans $W$) d'une
droite cyclique $D \subset \mathfrak{t}$ \uncertain{est} \ill{}
\struck{\ill{}} \uncertain{vectorielle} … \uncertain{et} $D/\mu_h$ a une
structure vectorielle \struck{\ill{}}. Il serait intéressant de calculer
\struck{la relation de \ill{}} la \uncertain{valeur} \ill{} de $G$
(\uncertain{connexe} \uncertain{cal.}) dans $\mathbb{Z}/2\mathbb{Z} \simeq
\mu_2(\mathbb{Z})$ \uncertain{qui} \uncertain{s'obtient} ainsi en \ill{}
\uncertain{et} le \uncertain{can.} \uncertain{provenant} de la \uncertain{vraie}
\ill{} et \ill{} \uncertain{compare} avec le car.\ \uncertain{privilégié}
\ill{} \uncertain{racines}). \uncertain{P.\ ex.}\ il est clair \ill{}
\ill{} $\mathfrak{g}_\psi$ ($\psi$ plus grande racine). \ill{} \ill{}
résultats de Kostant que si la deuxième \uncertain{racine} \ill{} \ill{}
\ill{} \ill{} \ill{} \ill{}, \uncertain{par} le tore $T$ dans un Borel $B$,
\ill{}, la primitive \uncertain{associée}. \uncertain{Alors} \ill{}
\[
  \mathfrak{g}_\psi \simeq \Delta \quad (\text{ou } \mathfrak{g}_\psi \simeq
  \Delta^{-1}, \text{ car } \Delta^{\otimes 2} \simeq \mathbf{1}),
\]
\ill{} \uncertain{isom.\ can.}\ \ill{} \ill{} \uncertain{isom.\ can.}
\note{les trois dernières lignes sont d'une écriture rapide et effacée ; on
n'en donne que ce qui se lit. En dessous, à l'encre brune :} \uncertain{Vérifier}
\ill{} $=$, \ill{} …

\marginal{\emph{NB} Si on a plus \uncertain{généralement} un topos, \ill{}
classifiant $B_{E_k^{*}}$, on \uncertain{trouve} \ill{} \uncertain{généralement}
\ill{} \ill{} \ill{}, donc on a un \uncertain{torseur}, \uncertain{défini} \ill{}
Kost Kill norm : \ill{} \ill{} \uncertain{faut} \ill{} \ill{} $\Delta$ \ill{}
\ill{}. Il faut que \uncertain{cela} provienne \ill{} le \uncertain{torseur}
\ill{} $E$ \ill{} $\Delta$ \ill{} être \uncertain{un} trivial. Il faut
\ill{} \uncertain{processus} à \uncertain{vérifier} \ill{} d'\ill{} et
\ill{} \struck{\ill{}} \uncertain{des} \emph{façon} générale \ill{} \ill{}
\ill{} obstruction \struck{\ill{}} \uncertain{de} Kostant $D$ dans
$H^{2}(Q ; E ; \mu_h)$, \uncertain{\ill{}} \uncertain{nul} \uncertain{si}
$H^{2}(Q, \mu_h) = \ill{}$ \ill{} ; \uncertain{restreint} :
$H^{2}(Q ; E^{\circ} ; \mu_h)$}
\note{écrit en oblique dans la marge gauche, du bas vers le haut, en regard
de b) ; les deux $H^2$ et « restreint » sont à l'encre brune. La lecture est
très lacunaire}

\page{34}
Considérons le schéma des \struck{\ill{}} \uncertain{rigidifications}
RigKos
\[
  (T, T', B, B') \quad \text{avec} \quad
  \begin{cases}
    T, T' \text{ tores maximaux en apposition} \\
    B, B' \text{ des Borels contenant } T, T' \text{ respectivement}
  \end{cases}
  \qquad [\,(h,p) = 1\,]
\]
(Rigidification de Kostant). C'est un schéma lisse sur $S$, sur lequel $G$
opère librement.
\note{« Rigidification de Kostant » est ajouté en marge gauche et souligné ;
« $(h,p)=1$ » est encadré à droite}

RigKos$/G \simeq$ schéma fini étale sur $S$, \struck{étale} de
\uncertain{rang} $W^{2}/\varphi(h)hz$, où
\[
  \begin{cases}
    W \text{ ordre du groupe de Weyl} \\
    h \text{ ordre de l'élt de Coxeter} \\
    \varphi \text{ l'indicatrice d'Euler} \\
    z \text{ ordre du centre}
  \end{cases}
\]

Il revient au même : une représentation du groupe \uncertain{$\Gamma$} des
automorphismes extérieurs dans un schéma fini \uncertain{étale} sur
$\mathrm{Spec}\,\mathbb{Z}[1/h]$ [le schéma des \emph{classes} de
rigidifications de Kostant de $G$ sur $\mathrm{Spec}(\mathbb{Z}[1/h])$]. Le
dernier schéma est \struck{\ill{}} \uncertain{il} \struck{de}
\uncertain{déployé} : prenons un \uncertain{épinglage}, \uncertain{définissant}
$\mu_h \xrightarrow{\kappa} T$ canonique, \uncertain{défini par épinglage},
$D$ engendré par
\[
  e_{\alpha_1} + \cdots + e_{\alpha_r} + e_{-\psi} = e
\]
(\uncertain{défini au signe près}).
\note{« canonique » et « défini par épinglage » sont écrits sous
$\mu_h \xrightarrow{\kappa} T$, « défini au signe près » sous $e_{-\psi}$}

On trouve $T'$ comme centralisateur \struck{\ill{}} de cet élément.
\uncertain{On a} \struck{\ill{}} Alors le schéma des Borels contenant $T$,
\uncertain{isomorphe} à $W_{\mathbb{Z}}$, et le schéma $Q'$ des Borels
contenant $T'$, qui est \uncertain{très} \uncertain{tordu}, \struck{\ill{}}
c'est un torseur sous $W' = N(T')/T'$. On a (dans $W \times W'$) le sous-groupe
$\Phi = N \cap N'$ extension de $(\mathbb{Z}/h\mathbb{Z})^{*}$ par
\struck{\ill{}} $\mathfrak{z} \times \mu_h$, où $\mathfrak{z} \subset W$
(centralisateur de $\kappa$) et $\mu_h \subset W'$. Il opère
(\uncertain{librement}) sur $Q \times Q'$ : torseur sous $W \times W'$, et on
forme
\[
  (Q \times Q')/\Phi .
\]
Le groupe des automorphismes extérieurs de $G$ y opère via son action par
$\pm 1$ sur $\mathfrak{g}_{-\psi}$, donc en fait, \struck{\ill{}} \uncertain{à
travers} l'opération \uncertain{sur} \ill{}. \uncertain{Si} $E$ est trivial, on
\uncertain{a} \uncertain{à} \uncertain{faire} pour \uncertain{un} autre
\uncertain{cas}.
\note{le schéma $Q$ des Borels contenant $T$ n'est pas nommé sur la page ;
$\Phi$ rend la lettre écrite au-dessus de « $N \cap N'$ », surchargée. Les
quatre dernières lignes sont à l'encre brune}

\page{35}
\note{Même tapuscrit, feuillet paginé « 6 », fait de bandes de papier collées,
en partie à l'encre bleue de la machine : « 2. Accouplements définis par une
biextension », n°~2.1 — pour $P$, $Q$ deux Groupes abéliens sur $T$ et
$n > 0$, construction d'un homomorphisme canonique (2.1.1)
$\alpha = \alpha^n_{P,Q} \colon {}_nP \otimes {}_nQ \to \mathrm{Tor}_1(P,Q)$,
composé (2.1.2) ; réduction au cas $nP = nQ = 0$ ; isomorphisme (2.1.3) dans
la catégorie dérivée, suite spectrale de Künneth (2.1.4), valeurs des
$\mathrm{Tor}_i^{\mathbb{Z}}(P, \mathbb{Z}/n\mathbb{Z})$ par la résolution
$0 \to \mathbb{Z}_T \xrightarrow{n} \mathbb{Z}_T \to 0$, suite exacte (2.1.5).
Sans rapport avec les notes sur Kostant. Corrections à la main : les $\alpha$
et le « $n$ » de la résolution encrés ; en marge, entouré, « $nQ = 0$. » et
une rature ; en marge droite, encadré, « (cf.\ SGA~1~IV~5.2) » ;
« de Künneth » ajouté après « suite spectrale de type homologique » ;
« qui figure dans la dernière suite exacte » précédé, au-dessus de
« aussi » biffé, de « l'homomorphisme-coin
$E^2_{0,1} \to \mathrm{Tor}_1^{\mathbb{Z}}$ » ; la dernière formule,
${}_nP \otimes Q \to \mathrm{Tor}_1({}_nP, Q) \to \mathrm{Tor}_1(P,Q)$, est
achevée à la main}

\section*{Liste d'espaces homogènes remarquables et leurs relations}

\note{titre de sa main, en haut à droite d'un feuillet de papier bis qui ne
porte rien d'autre (page 36) ; une flèche au crayon le prolonge par « liés
aux ss-groupes simples principaux de rang 1, ». Le feuillet est une chemise et
ne reçoit donc pas de \texttt{page}}

\page{37}
\textbf{Schémas :}

\begin{itemize}
  \item des sous-groupes simples (\uncertain{de rang 1}) principaux : SSP
  \item ---, avec épinglage : SSPEP $\subset
  \underline{\mathrm{Hom}}^{*}_{\mathrm{gr}}(\mathrm{Sl}(2), G)$
  \struck{\ill{} \ill{}}
  \item ---, avec couple de Kill : SSPKill
  \item ---, avec tore max. : SSPTor
  \item ---, avec tore maximal « \uncertain{généré} » de $\mathbb{G}_m$ :
  SSPTor$_0$ $\subset \underline{\mathrm{Hom}}_{\mathrm{gr}}(\mathbb{G}_m, G)$
  \item ---, avec Borel : SSPBor
  \item ---, avec Borel épinglé : SSPBorep
\end{itemize}

\begin{itemize}
  \item des tores (\struck{\ill{}} de dim.~1) principaux : STP
  \item des hom.\ $\mathbb{G}_m \to G$ principaux (cf.\ sous-tores
  « épinglés ») : ST$_0$P
\end{itemize}

\begin{itemize}
  \item des sous-groupes additifs principaux : SAP
  \item des sous-groupes additifs principaux pointés : SAP$_0$ $\subset
  \underline{\mathrm{Hom}}_{\mathrm{gr}}(\mathbb{G}_a, G)$
  \item des \struck{éléments} \ill{} unip.\ tout principaux : EUP $\subset G$
\end{itemize}
\marginal{sont isomorphes canoniquement car.~$0$ (et sans doute d'autres)}
\note{accolade en face des deux dernières lignes}

\struck{des éléments} \emph{NB} \ill{} \uncertain{peut} donner des variantes en
\uncertain{termes} de $\mathfrak{g}$ des trois \uncertain{dernières},
\uncertain{qui} sont \uncertain{isom.} \uncertain{car} \uncertain{les}
applications \uncertain{canoniques} \uncertain{isom.} en car.~$0$ \uncertain{ne}
\ill{} \ill{}.
\note{une tache d'encre couvre un mot après « NB » ; la lecture de cette
phrase est très incertaine}

\note{la suite du feuillet est coupée par trois longs traits obliques au
crayon qui la traversent ; ils ne biffent rien de lisible}

\begin{itemize}
  \item des éléments \struck{\ill{}} principaux \uncertain{en tous cas} (\ill{}
  \uncertain{réguliers}) : \struck{\ill{}} E Cox, \struck{\ill{}} E Kost
  \marginal{deviennent isomorphes une fois choisie une racine primitive
  $h$-ième de $1$ (\uncertain{donc} dans $\mu_h^{*}$ !)}
  \struck{(\emph{NB} \ill{} \ill{} il vaut mieux remplacer par les tores
  principaux $\mu_h \to G$ ?)}
\end{itemize}

\begin{itemize}
  \item des droites cycliques $\subset \mathfrak{g}$ : DC
  \item des \struck{\ill{}} éléments cycliques $\in \mathfrak{g}$ :
  DC$_0$ $\simeq$ EC
  \item des éléments \struck{principaux} cycliques principaux : ECP
  \struck{\ill{}} \struck{\ill{}} k$_0$
\end{itemize}
\marginal{\emph{NB} \uncertain{on a} \uncertain{un} \uncertain{espace}
homogène $G$ ; \ill{} DC$_0/G$ \struck{\ill{}} E$'$ \ill{} l'action \ill{}
les droites cycliques, l'action \ill{} \uncertain{normalisateur} \ill{}
$\mu_h$ \ill{} \ill{} \uncertain{multipliée} \ill{} \ill{} \ill{} $D/\Gamma$}
\note{encadré, à droite de ces trois lignes, en partie à l'encre brune ; lecture
très lacunaire}

\begin{itemize}
  \item des couples \struck{\ill{}} \struck{Coxeter} $(T, \gamma \in N(T))$ :
  \struck{\ill{}} \struck{\ill{}} \struck{Clox}
  \item des couples de tores maximaux en apposition : CTAP
  \item des couples de Coxeter $(T, \gamma \in W(T))$ : CCox
  \item des couples de Kostant $(\kappa : \mu_h \to G, D)$ : CKost
\end{itemize}
\marginal{sont isom.\ sur $\mu_h^{*}$}
\note{accolade en face des deux dernières lignes ; une flèche ramène la
deuxième ligne au-dessus de la première}

\page{38}
\note{sur le verso d'une page dactylographiée en français, qui transparaît à
l'envers et n'a pas de rapport avec ces notes. En haut, le diagramme des
espaces de la page précédente ; à droite, de sa main : « diagramme d'espaces
homogènes sous $G$ ». Les traits sont sans flèche}

\begin{tikzcd}[column sep=small, row sep=small, nodes={font=\scriptsize}]
  & & \underline{\mathrm{Isom}}(G_0, G) \arrow[d, no head, "\wr"] \arrow[r, dashed] & \underline{\mathrm{Hom}}_{\mathrm{gr}}(\mathrm{Sl}(2), G) & \\
  & & \text{SSPEP} \arrow[dl, no head] \arrow[d, no head] \arrow[dr, no head] & & \\
  & \text{SSPTor}_0 \arrow[r, no head, "\simeq"] \arrow[d, no head] \arrow[dl, no head] & \text{SSPKill} \arrow[dl, no head] \arrow[dr, no head] & \text{SSPBorep} \arrow[d, no head] \arrow[r, no head] & \text{SAP}_0 \simeq \text{EUP} \arrow[d, no head] \\
  \text{Kill} \simeq \text{ST}_0\text{P} \arrow[d, no head] & \text{SSPTor} \arrow[dr, no head] \arrow[dl, no head] & & \text{SSPBor} \arrow[dl, no head] \arrow[r, no head] & \text{SAP} \\
  \text{STP} & & \text{SSP} & &
\end{tikzcd}

\note{sous ST$_0$P, un trait épais descend vers « Tor » en bas à gauche ; à
côté, « Kost » surchargé d'un mot noirci, illisible. Plus bas, un second
diagramme, de groupes cette fois (les stabilisateurs, semble-t-il), les
étiquettes des traits étant les dimensions ou les quotients :}

\begin{tikzcd}[column sep=small, row sep=small, nodes={font=\scriptsize}]
  & & & e \arrow[dl, no head, "1"'] \arrow[d, no head, "1"] \arrow[dr, no head, "1"] & \\
  & & t \arrow[dl, no head, "r-1"'] \arrow[d, no head] \arrow[r, no head, "="] & t \arrow[dl, no head] \arrow[dr, no head, "1"] & b_u \arrow[d, no head, "1"] \arrow[dr, no head, "r-1"] & \\
  N_1 \arrow[r, no head, "\circ"] \arrow[d, no head] & T \arrow[d, no head, "\mathbb{Z}/2\mathbb{Z}"] & n \arrow[dr, no head, "2"'] \arrow[dl, no head, "r-1"'] & & b \arrow[dl, no head, "1"] \arrow[dr, no head, "r-1"] & Z(b_u) \arrow[d, no head, "1(\mathbb{G}_m)"] \\
  N \arrow[r, no head] & N_0 & & H & & N(b_u)
\end{tikzcd}

\note{un « ? » devant $N_1$ ; l'étiquette du trait vertical de $t$ à $n$ est
noircie et illisible. La disposition est simplifiée : sur la page, $t$ de
gauche est plus haut que $t$ de droite, et $N_1$, $T$, $N_0$, $N$ forment un
losange à gauche}


\page{39}
\note{Feuillet de tapuscrit du même exposé sur les topos que la page 27 (fin
de 3.1.2 et n°~3.1.3) : $u_*$
respecte toute structure algébrique définissable exclusivement en termes de
limites projectives, mais n'est en général pas exact à droite — source des
propriétés cohomologiques étudiées dans l'exposé suivant ; 3.1.3 : devant un
foncteur $f \colon E \to F$ d'un $\mathcal{U}$-topos dans un autre, expliciter
ses propriétés d'exactitude ; s'il commute aux limites inductives quelconques
et aux limites projectives finies, l'écrire $f = u^*$ pour un morphisme de
topos $u \colon F \to E$. Sans rapport avec les notes sur Kostant. Ses ajouts à
l'encre bleue :}

\marginal{les opérations fonctionnelles habituelles faites en termes de telles
structures, plus précisément à toutes les opérations qui \struck{\ill{}}
peuvent s'exprimer en termes de $\varinjlim$ et de $\varprojlim$ finies :
constructions d'objets libres (groupes ou modules libres, p.\ ex.) engendrés
par un objet, produits tensoriels (cf.\ §~6 plus bas) etc.}
\note{en tête du feuillet, encadré, relié par un trait à l'insertion
suivante}

\marginal{De plus, le foncteur $u^*$ « commute » à toutes}
\note{inséré après « limites inductives) », avec l'indication en marge
gauche : « à la ligne, puis : Quant au … » ; « Quant au foncteur image
directe $u_*$ » reçoit « $: E \to E'$ », « respecte » reçoit « par suite »,
« définissable en termes de limites projectives » reçoit « exclusivement »,
et « parmi » remplace un mot biffé}

\marginal{Par suite, il ne s'étend pas, en général, en présence d'}
\note{après « dans l'exposé suivant », « s'étend pas » au-dessus d'un mot
biffé ; appel d'une note de bas de page, de sa main :}

\marginal{ou un foncteur sur des objets du type comodule, ou bigèbre, ou
torseur sous un groupe et ne commute pas en général à des opérations telles
que « module libre engendré », produit tensoriel de modules, etc.}

\page{40}
\textbf{Résultats de Kostant} sur classes de conjugaison remarquables dans $G$
\uncertain{$\frac12$ simple} adjoint (sur $\mathbb{C}$)

a) Classes de conjugaison de ss-groupes simples de rang~$1$
\struck{\ill{}} [ou de ss-groupes additifs de $G$] $\simeq$ classes de
conjugaison d'éléments unipotents de $G$ [ou de $\mathfrak{g}$]
$\simeq$ certaines classes de conjugaison
\struck{d'éléments semi-simples de $\mathfrak{g}$ [\emph{NB} \ill{} \ill{}]}
d'homomorphismes $\mathbb{G}_m \to G$ [ou d'éléments $\in \mathfrak{g}$ … ].

Les hom.\ admissibles $\mathbb{G}_m \to T \simeq \mathbb{G}_m^{\Sigma}$
($\Sigma$ ensemble des racines simples) dans le « domaine fondamental » $W$ de
$\mathrm{Hom}(\mathbb{G}_m, T) \simeq \mathbb{Z}^{\Sigma}$, savoir
$\mathbb{N}^{\Sigma}$, correspondant à \emph{des} éléments dont les composantes
sont $0, 1, 2$ (pas tous).

b) Le cas où toutes les composantes sont $1$ est admissible ; il correspond
aux éléments unipotents principaux (\uncertain{ou} réguliers), et dans
$\mathfrak{g}$ aux éléments semi-simples principaux ;
\struck{\ill{}} aux hom.\ « principaux » $\mathbb{G}_m \to G$, de
\uncertain{part} des sous-groupes simples de rang~$1$ principaux.
\note{une grande marque en marge gauche, en face de b), illisible}

Centralisateur d'un élément unipotent principal : sous-groupe commutatif
(\uncertain{de Springer}) unipotent de dim.\ $r$ ($\simeq \mathbb{G}_m$
\uncertain{comme} \ill{} \uncertain{par} \uncertain{Kostant}).
\note{les mots « commutatif », « de Springer », « unipotent » et « de dim. »
sont des ajouts entourés au-dessus de la ligne ; leur ordre est incertain}

Normalisateur d'un ss-gr.\ additif principal : groupe \uncertain{connexe}
\struck{\ill{}} \ill{} \ill{}.

Centralisateur d'un hom.\ principal $\mathbb{G}_m \to G$ : un tore max.
\uncertain{T}. \uncertain{Donc} \ill{} bijection \uncertain{entre} les hom.\ et
les couples de Killing de $G$ \uncertain{correspondant} à~$b$.
\note{« Killing de $G$ » est ajouté au-dessus de la ligne}

Normalisateur d'un 1-tore principal : groupe \struck{\ill{}} des $g \in T$
\uncertain{d'ordre}~$2$ \struck{\ill{}} du rang~$1$ principal : lui-même.

\struck{Normalisateur} Stabilisateur d'un ss-groupe principal de rang~$1$
(avec un Borel, resp.\ un couple de Killing) : \uncertain{donné}
\uncertain{dans} \uncertain{chaque} \ill{} (resp.\ \ill{}) \ill{}
$H$ avec un tore maximal $T_0$ ; Normalisateur $N_0$ du tore $T_0$ dans $H$
(resp.\ $B_0$, resp.\ $T_0$).

\marginal{Donc sur base quelconque, $\exists$ ss-groupe \struck{\ill{}} de
rang~$1$ principal déployé ssi $G$ quasi-déployé.}
\note{encadré dans la marge gauche, relié par une flèche à « Normalisateur
d'un 1-tore principal »}

\marginal{\struck{\emph{NB} \ill{} \ill{} \ill{} \ill{} \ill{} \ill{} \ill{}
\ill{} \ill{}} \uncertain{question} \ill{} \uncertain{classes} \ill{} \ill{}
tores de dim.~$1$ principaux}
\note{encadré et biffé, sous le précédent}

\marginal{l'élément $w_0 \in W$ qui transforme la chambre $C$ en
l'opposée}
\note{écrit en long dans la marge gauche, en face de a) ; auquel des énoncés
il se rattache, la page ne le dit pas}

c) Soit $\mathbb{G}_m \xrightarrow{u} G$ un hom.\ principal, et soit
\struck{\ill{}} $h$ \uncertain{le nombre} de
\ill{} [$=\sigma(\psi)+1$] (la plus haute racine $\psi$, \emph{plus}~$1$).
[Coxeter ($=$ somme des coeff.\ de la plus haute racine $+1$), et
considérons $u(\zeta)$.] Soit $\zeta$ \struck{une} racine primitive $h$-ième
de~$1$, et considérons aussi \uncertain{pour} \uncertain{$\zeta$} \ill{}
\uncertain{variable}) ; les $u(\zeta)$ ($\zeta$ variable). On prouve que les
éléments obtenus sont \emph{réguliers} (\emph{principaux}), \uncertain{tous}
conjugués. Les éléments obtenus sont \uncertain{réguliers}, et
$\gamma \in W = N(T)/T$ \struck{\ill{}} de $G$. Si $T$ est un tore maximal,
et $\gamma \in W = N(T)/T$ un élément de Coxeter, on prouve que les
représentants de $\gamma$ \struck{\ill{}} sont \uncertain{les}
éléments \ill{} tores principaux (donc réguliers).
\note{le paragraphe c) est surchargé d'ajouts interlinéaires dont l'ordre ne se
reconstitue pas avec certitude ; on donne les segments dans l'ordre où ils se
lisent}

\marginal{\uncertain{ordre} qui divise $h$, il vaut mieux considérer la
classe d'hom.\ $\mu_h \to G$}
\note{encadré dans la marge gauche, en face de c)}

Centralisateur d'un élément régulier principal $\gamma$ : ss-groupe de
$N(T)$ \struck{\ill{}} \uncertain{extension} du ss-groupe de $W = W(T)$ qui
stabilise $\gamma$. (?) C'est probablement $T$, i.e.
\struck{\ill{}} le \uncertain{donné} de $\mu : \mathbb{G}_m \to G$ \uncertain{et}
$\gamma$ \uncertain{variant} : celle de $\mu : \mathbb{G}_m \to G$.
\note{le « ? » est entouré}

d) Considérons l'opération de $W$ sur
$\mathfrak{t} = \mathrm{Lie}(T)$. Il existe des \struck{\ill{}}
droites $D$ \uncertain{\ill{}} \struck{dans} $\gamma$
(\uncertain{op.\ de Coxeter}) dont les valeurs propres \uncertain{sont} une
racine primitive $h$-ième de~$1$, \uncertain{vérifiant} \uncertain{qu'elles}
\uncertain{qui sont} « régulières », et ces droites sont conjuguées sous
l'action de $W$ (donc : \uncertain{forment} \uncertain{un conjugaison}
\ill{})
\marginal{$G$ simple}
\note{dans la marge gauche, en face de d) ; la page s'arrête sur cette ligne}
\end{document}
